\section{Institutions, Intellectuals, Mystics}

\begin{quotex}
If we put the subjective interest of the Indian mind along with its tendency to arrive at a synthetic vision, we shall see how \textbf{monistic idealism} becomes the truth of things. To it the whole growth of Vedic thought points; on it are based the Buddhistic and the Brahmanical religions; it is the highest truth revealed to India. \flright{\textsc{Sarvepalli Radhakrishnan}, \emph{Indian Philosophy}, Vol 1}

Catholicism is truly the most perfect religion, just as modern European philosophy [i.e., absolute idealism] is the most perfect philosophy: together they are the highest creations of the Aryan spirit. \flright{\textsc{Giovanni Gentile}, \emph{On Religion}}

The prerogative of the human state is objectivity; the essential content of which is the \textbf{Absolute}. There is no knowledge without objectivity of the intelligence; there is no freedom without objectivity of the will; and there is no nobility without objectivity of the soul … Esoterism seeks to realize pure and direct objectvity; this is its raison d'être. \flright{\textsc{Frithjof Schuon}, \emph{Esoterism As Principle and As Way}}

\end{quotex}
\textbf{Friedrich von Hügel}, Baron of the Holy Roman Empire, identified in his book \emph{The Mystical Element of Religion} three elements of religion: \textbf{the institutional, the intellectual, and the mystical}. Although his intent should be clear from the names, I prefer to use more neutral names that can be defined with a specificity unencumbered by expectations, viz., the exoteric, mesoteric, and esoteric.

Historically, the three elements have been at odds with each other.

The doctrine of the three souls is clearly taught, even in exoteric theology. What is not so clearly taught is that each person has by birth his center of gravity in one of those souls. Thus a man, by nature, is centered on sensations, emotions or thoughts. No one of them is ``better" or ``higher" than the others, since each is the result of an interior imbalance.

Hence, the spiritual paths preferred by each type of person will differ. We can presume that a person centered on this-worldly sensations will be unlikely to be attracted to spirituality. In the cases where they do, the spiritual path is often of a low-grade, focusing on paranormal phenomena, séances, talismans, divination, or techniques promising to deliver money, sex, or power.

The emotion-centered type will gravitate to the exoteric element, and the thinking person to the mesoteric. The esoteric, on the other hand, is open only to those who manage to harmonize the three centers into a greater whole.

\paragraph{Exoteric}
Obviously the institutional church attracts adherents from all three types, but we will focus on the emotional types since they are least likely to move beyond it. For them, the emotional impact is the main focus. The mass and sacraments will make them feel good. They are attracted to devotional practices of a sentimental nature.

In extreme cases, as is not uncommon today, converts and reverts become emotionally involved in ``saving" the church, creating an Internet presence. These are the super-correct, even to the point of criticizing popes and bishops. At its highest, the emotional center evokes feelings of loyalty, solidarity, and fidelity. At its worse, these feelings combine with the incensive power to lash out everywhere. They are always hunting for heretics even if their own understanding is itself heretical and untraditional.

Unable to fully grasp the intellectual level, they regard theological ideas as dogmas. Hence, they interpret philosophical ideas in emotional terms, rather than as reasoned-out ideas. This leads to suspicion of those who restate metaphysical ideas in unaccustomed terms.

Similarly, mystics are viewed as just more intense versions of themselves. They develop emotional attachments to devotions of their favorite saints. Ultimately, the value of the saints and mystics lies in their own adherence to the exoteric, institutional church.

Scriptures are understood as naturalistic, historical, and scientific texts. That is a dangerous position to hold, because it justifiably permits historians and scientists to apply their own methods to those texts. In that atmosphere, the commitment to the moral, allegorical, and mystical interpretation of the texts is quite weak.

This is clearly a one-sided representation, but the purpose is only to point out those misconceptions that tend to create frictions and misunderstandings between the institutional and other perspectives. Although there is the danger of being cast out as heretical, esoterists need to respect exoteric practices. This is the teaching both of a \textbf{Rene Guenon} and a \textbf{Valentin Tomberg}, who insist on an accommodation between the exoteric and the esoteric.

\paragraph{Mesoteric}
The ancients certainly recognized that \textbf{Plato}, \textbf{Hermes Trismegistus}, and others held a true conception of God. It is a dogma of the Catholic Church that the God known to philosophy is the same God as that of revelation. There are several ways to approach this. Most often, theologians and philosophers arise from the exoteric church, and reconsider its teachings and dogmas in rational and intellectual terms.

In other cases, philosophers have developed conceptions of God apart from the exoteric church. As the epigraph from Radhakrishnan shows, these take the form of some sort of monist or absolute idealism at their best. At their worst, they promote a form of theistic personalism\footnote{\url{https://edwardfeser.blogspot.com/2013/04/craig-on-theistic-personalism.html}}, i.e., the conception that God is one being among others, albeit of a very special kind.

Now it is dogmatic that God is both Absolute and Infinite. These concepts are certainly comprehensible to the rational mind. It is well worth the effort to work through the demonstrations in such works as \emph{Multiple States of Being} by \textbf{Rene Guenon} or \emph{Survey of Metaphysics and Esoterism} by \textbf{Frithjof Schuon}. Theistic personalism denies both those attributes.

Radhakrishnan, in his massive study of Indian philosophy, gives many examples of Western philosophers whose thinking aligns with the Vedanta. Inter alios, he writes about \textbf{Francis Bradley}, who represents the peak of British idealism, \textbf{Giovanni Gentile}, who does the same for Italian idealism, and by extension, for continental European idealism, and even Baron von Huegel.

Since thought is not self-sufficient, i.e., it first needs something to think about, when the intellectual types divorce themselves from the exoteric element, it becomes free-floating. The tendency is to make rather simple things more complex than they are. Hence, they falsely regard their philosophy as superior to the exoteric. The result, all too often, is that they align themselves with the consensus worldview of conventionally educated people. Hence, their erudition and skills in logic and argument produce no more than what the average person can come up with.

At its worse, the intellectual takes uncertainty and doubt as a ``badge of honor", a mark of superior intelligence. He will glibly talk about ``shades of grey", a fitting phrase for intellectual bondage, ``seeing both sides" of an issue, and so on. On the contrary, and this should be obvious to everyone, it is an admission that his intellectual reasoning cannot come to any sort of certainty.

The man who comes to know himself sees that it is in the nature of the intellect to have doubt, the ``yes" turns into a ``no". Someone undeveloped will stay in this step. On the other hand, those who seek to become balanced will search for a higher path, one that transcends the intellect.

\paragraph{Esoteric}
Of course, we don't want to completely conflate the truly esoteric with the mystical as used by the Baron. The mystical encompasses many things and is not identical to the esoteric. The primary difference is that the mystic follows a passive path whereas the esoterist follows an active path. The esoterist notices the shortcomings in the institutional and intellectual approaches. However, rather than abandoning them entirely, he seeks to balance them and integrate them into a larger whole.

The certainty of faith of the exoterists becomes the certainty of gnosis. While not rejecting the rationality of the intellectuals, historians, and scientists, he becomes aware of a higher mind that is trans-rational, not irrational. This higher mind, the end point of theosis, transcends the duality of the discursive, thinking mind.



\flrightit{Posted on 2014-11-23 by Cologero }

\begin{center}* * *\end{center}

\begin{footnotesize}\begin{sffamily}



\texttt{David on 2014-11-24 at 13:05 said: }

Very good text, I will have to look at the von Huegel. 

That being said, what about non-dualism within this scheme, be it Christian or non-Christian (Advaita-Vâda)?


\hfill

\texttt{Boreas on 2014-11-24 at 14:56 said: }

There is a book that deals with the doctrine of non-dualism within Christian and perennialist context. Here's the link for the book:

\url{http://www.sophiaperennis.com/books/christianity/christianity-and-the-doctrine-of-non-dualism/}

Just for the record, I haven't read the book myself, but it came to my mind when you questioned the existence of a non-dualist position within Christianity.


\hfill

\texttt{David on 2014-11-24 at 16:57 said: }

@Boreas : This is indeed a book I was referencing. Might I add that Vladimir Lossky is the one that adapted the term for Christian use since the "Monk" that wrote the other book hinted that he took it from him.

@Fractal: I think it might exist. I am neither a saint nor a genius, so instead of explaining it myself, I will direct you to the book that Boreas linked. I would suggest in french since it is the original language and it have a preface by Jean Tourniac, an esoterist who worked on René Guénon. The book from the Monk is both linked with Tradition (use of Guénon most notably), and very orthodox (St Bernard, Nicholas de Cuze, St Thomas Aquinas, etc.). It received the autorization from its superior in his order (Cirstercian).


\hfill

\texttt{Wayne Ferguson on 2014-11-25 at 10:34 said: }

Trinitarian panentheism ? Christian Nonduality

Apropos of Christianity and Vedanta, I am interested to learn more about Paul Deussen. Quoting Letter 19 of ``Meditations on the Tarot":

\begin{quotex}
In relation to the theme of the mission of intelligence being the way to intuition, it is relevant to indicate the fact that the philosophical work of Immanuel Kant —which scorched the pretensions to certainty of autonomous intelligence with respect to metaphysical things by demonstrating the set limits of knowledge possible to autonomous intelligence —has had an effect comparable to that of wind, which extinguishes the weak fire and which revives the strong fire: the one becoming sceptics and the others becoming mystics. Kant put an end to the speculative metaphysics of autonomous intelligence and opened up the way to a mysticism which non-autonomous intelligence or ``practical reason" (praktische Vernunft) is capable of—Kant's ``practical reason" being intelligence united to the wisdom of moral nature, i.e. intuition. Indeed, several times I have had occasion to observe the fact that with time Kantians become mystics —to name, for example, the German philosopher Paul Deussen, the author of a synthesis of Kantianism, Platonism and the Vedanta (cf. Paul Deussen, The Elements of Metaphysics; trsl. C. M. Duff, London-New York, 1894).

\end{quotex}

\url{http://en.wikipedia.org/wiki/Paul\_Deussen}


\hfill

\texttt{Wayne Ferguson on 2014-11-25 at 10:46 said: }

p.s. sorry for the poorly formatted quote in the comment above. Also, the ``?" is what appeared in place of the ``approximately equal" sign that I tried to post. 

Here is a link to an online copy of the Deussen book that Tomberg mentions:

\url{https://archive.org/details/outlineofvedanta00deus}


\hfill


\end{sffamily}\end{footnotesize}
